%! Characteristics of Exponential Functions — Unit 6, Lesson 6.0 — Exit Ticket (Answer Key)
\documentclass[10pt]{article}
\usepackage{algebra2-article}
\usepackage{algebra2-key}   % pulls in -boxes; do NOT also load -boxes
\newcommand{\TallMath}[1]{$\displaystyle #1\rule[-1.4em]{0pt}{3.2em}$}
\newcommand{\lgrid}[4]{%
  \draw[step=1, linegray!40, very thin] (#1,#3) grid (#2,#4);
  \draw[->, semithick, charcoal] (#1,0)--(#2,0) node[below right, font=\scriptsize]{$x$};
  \draw[->, semithick, charcoal] (0,#3)--(0,#4) node[left, font=\scriptsize]{$y$};}
\newcommand{\expcurve}[4]{\draw[very thick, forest, domain=#3:#4, samples=120, smooth]
  plot (\x,{#1*exp((#2)*(\x))});}

\begin{document}

\pageheader{Unit 6, Lesson 6.0}{Exit Ticket --- Answer Key}
\namedateperiod

\vspace{0.08in}

No notes. The graph below is $f(x) = 8\left(\frac12\right)^{x}$. Use it for questions 1 and 2.
\begin{center}
\begin{tikzpicture}[scale=0.42]
  \lgrid{-1}{7}{-1}{10}
  \begin{scope}
    \clip (-1,-1) rectangle (7,10);
    \expcurve{8}{-0.693147}{-0.32}{7}
  \end{scope}
  \draw[navy, dashed, semithick] (-1,0)--(7,0);
  \node[navy, font=\scriptsize, above right] at (4.6,0) {$y=0$};
  \fill[charcoal] (0,8) circle (3pt) node[right, font=\scriptsize]{$(0,8)$};
  \fill[charcoal] (1,4) circle (3pt) node[right, font=\scriptsize]{$(1,4)$};
  \fill[charcoal] (2,2) circle (3pt) node[above right, font=\scriptsize]{$(2,2)$};
  \fill[charcoal] (3,1) circle (3pt) node[above right, font=\scriptsize]{$(3,1)$};
\end{tikzpicture}
\end{center}

\begin{enumerate}[leftmargin=1.8em, itemsep=8pt]
  \item \textbf{Read the features.} Write the domain and range in interval notation.

        $a=$ \ans{8} \quad $b=$ \ans{$\frac12$} \quad Growth or decay? \ans{decay}

        \vspace{3pt}
        Domain: \ans{$(-\infty,\infty)$} \quad Range: \ans{$(0,\infty)$} \quad
        Horizontal asymptote: \ans{$y=0$}

        \vspace{3pt}
        $y$-intercept: \ans{$(0,8)$} \quad $x$-intercept: \ans{none} \quad
        $f$ is \ans{decreasing} on its whole domain.

        \vspace{3pt}
        End behavior: as $x\to+\infty$, $f(x)\to$ \ans{$0$}; \; as $x\to-\infty$,
        $f(x)\to$ \ans{$+\infty$}.

  \item \textbf{Explain.} Why does $f$ have no $x$-intercept? And why does the graph have no
        absolute minimum, even though it clearly flattens out along the bottom?
        \ansline{$\left(\frac12\right)^{x}$ is positive for \emph{every} real $x$ --- a negative
        exponent produces a fraction, not a negative number --- so $8\left(\frac12\right)^{x}$ is
        always positive and never equals $0$. The graph never reaches the $x$-axis, so there is no
        $x$-intercept.\\[3pt]
        The flat part along the bottom is the \emph{horizontal asymptote} $y=0$. The graph gets as
        close to it as you like but never lands on it, and a minimum has to be a value the function
        actually reaches. So $0$ is a boundary of the range, not a minimum --- which is why the
        range is $(0,\infty)$ and not $[0,\infty)$.}

  \item \textbf{SOL-style multiple choice.} What is the range of $g(x) = 6\left(\frac13\right)^{x}$?
        \begin{enumerate}[label=\Alph*., leftmargin=1.9em, itemsep=1pt]
          \item $(0,\infty)$ \quad \textcolor{keyred}{\textbf{$\leftarrow$ correct}}
          \item $[0,\infty)$
          \item $(-\infty,\infty)$
          \item $(-\infty,0)$
        \end{enumerate}
\end{enumerate}

\begin{teachernote}
\textbf{Answer 3:} A. Every distractor breaks exactly one idea. \textbf{B} is \emph{the} signature
error of the lesson --- treating the asymptote as a value the function attains, exactly the Unit 5
endpoint habit carried over. \textbf{C} forgets the floor entirely and answers with the
\emph{domain}. \textbf{D} confuses ``decay'' with ``negative'': the outputs shrink, but they shrink
toward $0$ from above and never cross it.

\textbf{Sort the collected tickets into four piles} to decide how Lesson~6.1 opens:
\emph{got it} / \emph{closed the bracket} (item 3 answered B, or item 1's range written
$[0,\infty)$) / \emph{decay means negative} (item 3 answered D, or the wrong end-behavior arms in
item 1) / \emph{hunted for an $x$-intercept} (item 1 or 2). The first two piles are the same
misconception at different strengths and are worth a two-minute re-teach of the Section 4 picture.
Lesson 6.1 asks students to \emph{build} $y=ab^{x}$ from a story, so also check item 1's $a$ and $b$:
anyone who could not pull $a=8$ off the $y$-intercept will struggle immediately tomorrow.
\end{teachernote}

\end{document}
